\documentclass[aspectratio=169]{beamer}
\usetheme{simple}

\input{preamble/preamble.tex}
\input{preamble/preamble_math.tex}
% \input{preamble/preamble_acronyms.tex}

\title{Lecture 5: Sparse GLMs}
\subtitle{STATS305B: Applied  Statistics II}
\author{Scott Linderman}
\date{\today}


\begin{document}


\maketitle
    
\begin{frame}{Recap}

Last time...
\begin{itemize}
    \item Generalized Linear Models (GLMs)
    \item Canonical form
    \item Demo: Poisson GLM
    \item Non-canonical forms
    \item Model checking and comparison
\end{itemize}
    
\end{frame}
    
\begin{frame}{Outline}

    Today...
    \begin{itemize}
        \item (Weighted) LASSO Regression
        \item Synthetic Demo
        \item The GLMNet Algorithm
        \item Why it Works: Proximal Methods
    \end{itemize}
    
\end{frame}

    
\begin{frame}[t,allowframebreaks]{Motivation}\label{sparse-glms}

We like linear and generalized linear models because the parameters are readily interpretable. 

\(\beta_j\) relates changes in covariate \(x_j\) to changes in the natural parameter of the
response distribution. 

One common application of such models is for \emph{variable selection}, finding a subset of covariates that are most
predictive of the response. 

For variable selection, we would like our estimates, \(\hat{\mbbeta}\), to be \emph{sparse}. 

When we have a vast number covariates --- as in genome-wide association studies (GWAS) where we aim
to predict a trait given thousands of single nucleotide polymorphisms
(SNPs) in the genome --- sparse solutions help focus our attention on
the most relevant covariates.

\end{frame}

\begin{frame}[t,allowframebreaks]{Lasso Regression}\label{lasso-regression}

Consider a linear Gaussian model, 
\begin{align*}
y_i & \ind{\sim} \mathrm{N}(\beta_0 + \mbx_i^\top \mbbeta, 1) \qquad \text{for } i = 1, \ldots,n,
\end{align*} where \(\beta_0\) is the intercept
parameter, \(\mbx_i \in \reals^p\) are the covariates, and
\(\mbbeta \in \reals^p\) are the weights. We factored out the intercept
because we typically don't regularize that parameter.

The Lasso yields sparse solutions for linear models like this by
minimizing the (average) negative log likelihood subject to \(\ell_1\)
regularization. \begin{align*}
\cL(\beta_0, \mbbeta) 
&= -\frac{1}{n} \sum_{n=1}^n \log \mathrm{N}(y_i \mid \beta_0 + \mbx_i^\top \mbbeta, 1) + \lambda \|\mbbeta\|_1 \\
&= \frac{1}{2n} \sum_{i=1}^n (y_i - \beta_0 - \mbx_i^\top \mbbeta)^2 + \lambda \|\mbbeta\|_1 \\
&= \frac{1}{2n} \sum_{i=1}^n (y_i - \beta_0 - \mbx_i^\top \mbbeta)^2 + \lambda \sum_{j=1}^p |\beta_j|
\end{align*} 
This is still a convex objective function!

It's tempting to just use vanilla gradient ascent to find the optimum,
\begin{align*}
\mbbeta^{(t+1)} &\leftarrow \mbbeta^{(t)} - \alpha_t \nabla \cL(\mbbeta^{(t)}),
\end{align*} 

Unfortunately, the Lasso objective it is not continuously
differentiable: the \textbf{gradient at \(\beta_j=0\) is discontinuous}
due to the absolute value in the \(\ell_1\) norm. What can we do
instead? 
\begin{itemize}
    \item You only live once\ldots{} run gradient descent and hope for
the best! 
    \item Use \emph{subgradient} descent, taking a step in the
direction of any subgradient of \(\cL\), but that approach can be much
slower, with convergence rates of only \(\cO(1/\epsilon^2)\). 
    \item Use
\emph{proximal} gradient descent, which amounts to \emph{iterative soft
thresholding} for the Lasso problem. This has much better convergence
rates, as we'll discuss below. 
    \item Use \emph{coordinate descent}, which
works very well in practice, even if it's harder to prove convergence
rates.
\end{itemize}
\end{frame}

\begin{frame}[t,allowframebreaks]{Coordinate Descent}\label{coordinate-descent}

Fix all parameters except for \(\beta_j\) for some
\(j \in \{1,\ldots,p\}\). As a function of \(\beta_j\), the average
negative log likelihood is, \begin{align*}
\cL(\beta_j; \mbbeta_{\neg j}) 
&= \frac{1}{2n} \sum_{i=1}^n (\widetilde{y}_{ij} - x_{ij} \beta_j)^2 + \lambda |\beta_j| + c,
\end{align*} where \begin{align*}
\widetilde{y}_{ij} &= y_i - \beta_0 - \sum_{k \neq j} x_{ik} \beta_k,
\end{align*} and where \(c\) is constant with respect to \(\beta_j\).

This is just a scalar minimization problem. Completing the square, we
can rewrite the objective as, \begin{align*}
\cL(\beta_j; \mbbeta_{\neg j}) 
&= \frac{1}{2} \frac{(\beta_j - \mu_j)^2}{\sigma_j^2} + \lambda |\beta_j| + c'
\end{align*} 

\textbf{Exercise:} Solve for $\mu_j$ and $\sigma_j^2$.

\end{frame}

\begin{frame}[t,allowframebreaks]{Coordinate Descent}\label{coordinate-descent}
Let's plot the objective for a few values of $\mu_j$, fixing $\sigma_j^2 \lambda = 1$.
\begin{center}
\includegraphics[width=0.95\linewidth]{06_sparse_glms/output_6_0.png}
\end{center}
{ \hspace*{\fill} \\}
    
    Note that many of the minimizers (denoted by the red lines) are obtained
at \(\beta_j = 0\)! 
\end{frame}

\begin{frame}[t,allowframebreaks]{The Soft-Thresholding Operator}\label{soft-thresholding}

With a bit of calculus, we can show that the
minimizer is given by the \textbf{soft-thresholding operator},
\begin{align*}
\beta_j^\star 
&= \begin{cases}
\mu_j - \sigma_j^2 \lambda & \text{if } \mu_j > \sigma_j^2 \lambda \\
0 &\text{if } |\mu_j| < \sigma_j^2 \lambda \\
\mu_j + \sigma_j^2 \lambda &\text{if } u_j < -\sigma_j^2 \lambda
\end{cases} \\
&\triangleq S_{\sigma_j^2 \lambda}(\mu_j)
\end{align*}

We can write the soft-thresholding operator more compactly as,
\begin{align*}
S_\alpha(\mu) &= \mathrm{sign}(\mu) \max\{|\mu| - \alpha, 0\}.
\end{align*} 
\end{frame}

\begin{frame}[t,allowframebreaks]{Coordinate ascent step for the
intercept}\label{coordinate-ascent-step-for-the-intercept}

\textbf{Exercise:} Show that the coordinate update
for the intercept is, 
\begin{align*}
\beta_0^\star &= \frac{1}{n} \sum_{i=1}^n \widetilde{y}_{i0},
\end{align*} 
where \(\widetilde{y}_{i0} = y_i - \mbx_i^\top \mbbeta\).

\end{frame}

\begin{frame}[t,allowframebreaks]{Weighted Lasso
Regression}\label{weighted-lasso-regression}

Finally, suppose we have heteroskedastic noise, \begin{align*}
y_i &\sim \mathrm{N}(\beta_0 + \mbx_i^\top \mbbeta, w_i^{-1})
\end{align*} where \(w_i\) is the inverse variance (precision) of the
\(i\)-th observation.

Then the objective would become, \begin{align*}
\cL(\beta_0, \mbbeta) 
&= -\frac{1}{n} \sum_{n=1}^n \log \mathrm{N}(y_i \mid \beta_0 + \mbx_i^\top \mbbeta, w_i^{-1}) + \lambda \|\mbbeta\|_1 \\
&= \frac{1}{2n} \sum_{i=1}^n w_i (y_i - \beta_0 - \mbx_i^\top \mbbeta)^2 + \lambda \|\mbbeta\|_1
\end{align*} Note that the \(w_i\)'s become the \textbf{weights} in the
objective.

Following the same steps as above, we can write the objective as a
function of \(\beta_j\), \begin{align*}
\cL(\beta_j; \mbbeta_{\neg j}) 
&= \frac{1}{2} \frac{(\beta_j - \mu_j)^2}{\sigma_j^2} + \lambda |\beta_j| + c'
\end{align*} where \begin{align*}
\mu_j &= \frac{h_j}{J_j} \\
\sigma_j^2 &= \frac{1}{J_j} \\
J_j &= \frac{1}{n} \sum_{i=1}^n w_i x_{ij}^2 \\
h_j &= \frac{1}{n} \sum_{i=1}^n w_i \widetilde{y}_{ij} x_{ij}.
\end{align*} Again, the coordinate-wise minimum is obtained at
\(\beta_j^\star = S_{\sigma_j^2 \lambda}(\mu_j)\).

For the intercept,
\(\beta_0^\star = \frac{\sum_{i=1}^n w_i \widetilde{y}_{i0}}{\sum_{i=1}^n w_i}\).
\end{frame}

\begin{frame}[t,allowframebreaks]{Synthetic Demo}\label{synthetic-demo}

\begin{center}
    See Colab
\end{center}
\end{frame}

\begin{frame}[t,allowframebreaks]{Questions}\label{questions}
    
\begin{enumerate}
\def\labelenumi{\arabic{enumi}.}
\item
  Can you identify two was in which the Lasso estimate differs from the
  OLS estimate?
\item
  What estimate do you get if you set \(\lambda = 0\)?
\item
  What if you take \(\lambda \to \infty\)? 
\end{enumerate}
\end{frame}

\begin{frame}[t,allowframebreaks]{Fitting Sparse GLMs}\label{fitting-sparse-glms}

Now let's generalize this approach to fit \(\ell_1\)-regularized GLMs!
This is exactly what the \texttt{glmnet} package
(Friedmann et al., 2010) solves.

Suppose we have a GLM with the canonical mean function, \begin{align*}
\E[Y_i] &= f(\beta_0 + \mbx_i^\top \mbbeta)
\end{align*} where we have again factored out the intercept.

The regularized objective is, \begin{align*}
\cL(\beta_0, \mbbeta)
&= -\frac{1}{n} \sum_{i=1}^n \log p(y_i \mid f(\beta_0 + \mbx_i^\top \mbbeta)) + \lambda \|\mbbeta \|_1.
\end{align*}
\end{frame}

\begin{frame}[t,allowframebreaks]{Review: Iteratively Reweighted Least
Squares}\label{review-iteratively-reweighted-least-squares}

Recall that Newton's method for canonical GLMs (without regularization)
is equivalent to \textbf{iteratively reweighted least squares}. The
\((t+1)\)-th step of Newton's method is equivalent to solving a weighted
least squares problem to find the minimum of an objective,
\begin{align*}
\widetilde{\cL}(\beta_0, \mbbeta; \beta_0^{(t)}, \mbbeta^{(t)})
&= \frac{1}{n} \sum_{i=1}^n w_i^{(t)} (z_i^{(t)} - \beta_0 - \mbx_i^\top \mbbeta)^2
\end{align*} where the working responses are \begin{align*}
z_i^{(t)} = \beta_0^{(t)} + \mbx_i^\top \mbbeta^{(t)} + \frac{y_i - \hat{y}_i^{(t)}}{w_i^{(t)}},
\end{align*} the predictions are, \begin{align*}
\hat{y}_i^{(t)} &= f(\beta_0^{(t)} + \mbx_i^\top \mbbeta^{(t)}),
\end{align*} and the weights are equal to the conditional variances,
\begin{align*}
w_i^{(t)} &= \mathrm{Var}[Y_i \mid \mbx_i, \beta_0^{(t)}, \mbbeta^{(t)}].
\end{align*}

For example, in a logistic regression, \begin{align*}
\hat{y}_i^{(t)} &= \sigma(\beta_0^{(t)} + \mbx_i^\top \mbbeta^{(t)}) \\
w_i &= \hat{y}_i^{(t)} (1 - \hat{y}_i^{(t)}).
\end{align*}
\end{frame}

\begin{frame}[t,allowframebreaks]{The Algorithm}\label{the-algorithm}

From here, we can sketch out a pretty straightforward algorithm for
fitting sparse GLMs. 

Within each Newton iteration, solve a weighted
least squares problem, \emph{subject to the \(\ell_1\)-regularization
penalty}, using coordinate descent. 

Once the coordinate descent
procedure converges, update the working responses and weights, then
repeat. 

This is essentially the algorithm in \texttt{glmnet}
(Friedmann et al., 2010)!
\end{frame}

\begin{frame}[t,allowframebreaks]{Synthetic Demo}\label{synthetic-demo-2}

\begin{center}
    See Colab
\end{center}
\end{frame}

\begin{frame}[t,allowframebreaks]{Computational Tricks}\label{computational-tricks}

There are several simple tricks to speed it up. 

\begin{itemize}
    \item Rather than recomputing the residual for each
coordinate, we can update and downdate the residual after each
coordinate update. 

    \item You can show that the coordinate updates only
depend on sufficient statistics \(\sum_{i} w_i x_{ij} x_{ik}\) and
\(\sum_i w_i z_i x_{ij}\), and these statistics don't change within the
each outer loop. We can save some time by precomputing these at the
start of the \texttt{\_glmnet\_step} function. This trick is referred to
as using \emph{covariance updates}. 

    \item When \(X\) is sparse, we can
implement the sufficient statistics calculations even more efficiently.
\end{itemize}

Friedmann et al (2010) describe several other
implementation-level details for making the code as fast as possible.
You should also check out James Yang's thesis and his
\href{https://github.com/JamesYang007/adelie}{amazing \texttt{adelie}
package for LASSO problems}.
\end{frame}

\begin{frame}[t,allowframebreaks]{Proximal Methods}\label{proximal-methods}

The \texttt{glmnet} algorithm is intuitive, but \textbf{why the heck
does it work?!} To gain a deeper theoretical understanding, let's take a
step back and talk about proximal methods.
\end{frame}

\begin{frame}[t,allowframebreaks]{Proximal Gradient
Descent}\label{proximal-gradient-descent}

Proximal gradient descent is an optimization algorithm for convex
objectives that decompose into a differentiable part and a
non-differentiable part, \begin{align*}
\cL(\mbbeta) &= \cL_{\mathsf{d}}(\mbbeta) + \cL_{\mathsf{nd}}(\mbbeta)
\end{align*} where \(\cL_{\mathsf{d}}\) is convex and
\emph{differentiable}, whereas \(\cL_{\mathsf{nd}}\) is convex but
\emph{not differentiable}. The idea is to stick as close to vanilla
gradient descent as possible, while correcting for the
non-differentiable part of the objective.

If we just had the differentiable part, \(\cL_{\mathsf{d}}\), we could
perform gradient descent. One way to think about the gradient descent
update is as the solution to a quadratic minimization problem,
\begin{align*}
\mbbeta^{(t+1)} &\leftarrow 
\arg \min_{\mbz} \cL_{\mathsf{d}}(\mbbeta^{(t)}) + \nabla \cL_{\mathsf{d}}(\mbbeta^{(t)})^\top (\mbz - \mbbeta^{(t)}) + \frac{1}{2 \alpha_t} \| \mbz - \mbbeta^{(t)}\|_2^2
\end{align*} We can think of the surrogate problem as a second order
approximation of the objective in which the Hessian is replaced with
\(\frac{1}{\alpha_t} \mbI\).

\emph{Proximal} gradient descent follows the same logic, but it keeps
the non-differentiable part, \begin{align*}
\mbbeta^{(t+1)} &\leftarrow 
\arg \min_{\mbz} \cL_{\mathsf{d}}(\mbbeta^{(t)}) + \nabla \cL_{\mathsf{d}}(\mbbeta^{(t)})^\top (\mbz - \mbbeta^{(t)}) + \frac{1}{2 \alpha_t} \| \mbz - \mbbeta^{(t)}\|_2^2 + \cL_{\mathsf{nd}}(\mbbeta^{(t)}) \\
&= \arg \min_{\mbz} \frac{1}{2 \alpha_t} \| \mbz - (\mbbeta^{(t)} - \alpha_t \nabla \cL_{\mathsf{d}}(\mbbeta^{(t)})) \|_2^2 + \cL_{\mathsf{nd}}(\mbz) \\
\end{align*} The resulting update balances two parts: 
\begin{enumerate}
    \item  Stay close to
the vanilla gradient descent update,
\(\mbbeta^{(t)} - \alpha_t \nabla \cL_{\mathsf{d}}(\mbbeta^{(t)})\).

    \item Also minimize the non-differentiable part of the objective,
\(\cL_{\mathsf{nd}}(\mbbeta^{(t)})\).
\end{enumerate}

As a sanity check, note that we recover vanilla gradient descent with
\(\cL_{\mathsf{nd}}(\mbbeta^{(t)}) = 0\).
\end{frame}

\begin{frame}[t,allowframebreaks]{Proximal Mapping}\label{proximal-mapping}

We call the function, \begin{align}
\mathrm{prox}(\mbu; \alpha_t) 
&= \arg \min_{\mbz} \frac{1}{2 \alpha_t} \| \mbz - \mbu \|_2^2 + \cL_{\mathsf{nd}}(\mbz)
\end{align} the \textbf{proximal mapping}.

Notes
\begin{itemize}
    \item The proximal mapping depends on the form of the non-differentiable part of the objective, even though we
have suppressed that in the notation. 
    \item However, it does \emph{not} depend on the form of the continuous part of the objective. 
\end{itemize}
\end{frame}

\begin{frame}[t,allowframebreaks]{Algorithm}\label{algorithm}

With this definition, the proximal gradient descent algorithm is,

\textit{Proximal Gradient Descent}

\textbf{Input:} Initial parameters \(\mbbeta^{(0)}\), proximal mapping
\(\mathrm{prox}(\cdot; \cdot)\).

\begin{itemize}
\item
  \textbf{For} \(t=1,\ldots, T\)

  \begin{itemize}
  \item Set
    \(\mbbeta^{(t)} \leftarrow \mathrm{prox}(\mbbeta^{(t-1)} - \alpha_t \nabla \cL_{\mathsf{d}}(\mbbeta^{(t-1)}); \alpha_t)\).
  \end{itemize}
\end{itemize}

\textbf{Return} \(\mbbeta^{(T)}\). 

So far, it's not obvious that this framing is helpful. We still have a
potentially challenging optimization problem to solve in computing the
proximal mapping. However, for many problems of interest, the proximal
mapping has simpled closed solutions.
\end{frame}

\begin{frame}[t,allowframebreaks]{Proximal Gradient Descent for Lasso
Regression}\label{proximal-gradient-descent-for-lasso-regression}

Consider the Lasso problem. The objective decomposes into convex
differentiable and non-differentiable parts, \begin{align*}
\cL_{\mathsf{d}}(\mbbeta) &= \frac{1}{2} \|\mby - \mbX \mbbeta\|_2^2 \\
\cL_{\mathsf{nd}}(\mbbeta) &= \lambda \|\mbbeta\|_1.
\end{align*}
\end{frame}

\begin{frame}[t,allowframebreaks]{Proximal Mapping}\label{proximal-mapping}

The proximal mapping is, \begin{align*}
\mathrm{prox}(\mbu; \alpha_t) 
&= \arg \min_{\mbz} \frac{1}{2 \alpha_t} \| \mbz - \mbu \|_2^2 + \lambda \|\mbz\|_1 \\
&= \arg \min_{\mbz} \sum_{j=1}^p  \frac{1}{2 \alpha_t} (z_j - u_j)^2 + \lambda |z_j| 
\end{align*} It separates into optimization problems for each
coordinate, and each coordinate has a closed-form solution in terms of
the soft-thresholding operator! \begin{align*}
[\mathrm{prox}(\mbu; \alpha_t)]_j 
&= S_{\alpha_t \lambda}(u_j)
\end{align*}
\end{frame}

\begin{frame}[t,allowframebreaks]{Iterative Soft-Thresholding
Algorithm}\label{iterative-soft-thresholding-algorithm}

Now let's plug in the gradient of the differentiable part,
\begin{align*}
\nabla \cL_{\mathsf{d}}(\mbbeta) &= \mbX^\top (\mby - \mbX \mbbeta).
\end{align*} Substituting this into the proximal gradient descent
algorithm yields what is sometimes called the \textbf{iterative
soft-thresholding algorithm (ISTA)},

\textit{Iterative Soft-Thresholding}

\textbf{Input:} Initial parameters \(\mbbeta^{(0)}\), covariates
\(\mbX \in \reals^{n \times p}\), responses \(\mby \in \reals^n\)

\begin{itemize}
\item
  \textbf{For} \(t=1,\ldots, T\)

  \begin{itemize}
  \item Set
    \(\mbbeta^{(t)} \leftarrow S_{\alpha_t \lambda}(\mbbeta^{(t-1)} - \alpha_t \mbX^\top (\mby - \mbX \mbbeta^{(t-1)}))\).
  \end{itemize}
\end{itemize}

\textbf{Return} \(\mbbeta^{(T)}\). 
\end{frame}

\begin{frame}[t,allowframebreaks]{Convergence}\label{convergence}

If \(\nabla \cL_{\mathsf{d}}\) is \(L\)-smooth then proximal gradient
descent with fixed step size \(\alpha_t = 1/L\) then, \begin{align*}
f(\mbbeta^{(t)}) - f(\mbbeta^\star) \leq \frac{L}{2 t} \|\mbbeta^{(0)} - \mbbeta^\star \|_2^2,
\end{align*} so it matches the gradient descent convergence rate of
\(\cO(1/\epsilon)\). (With Nesterov's accelerated gradient techniques,
you can speed this up to \(\cO(1/\sqrt{\epsilon})\).
\end{frame}

\begin{frame}[t,allowframebreaks]{Proximal Newton Method}\label{proximal-newton-method}

One great thing about proximal gradient descent is its generality. We
could easily apply it to \(\ell_1\)-regularized GLMs, substituting the
gradient of the negative log likelihood, which also has a simple closed
form expression. The proximal operator remains the same, and we obtain
the same converge rates as gradient descent on standard GLMs.

However, we saw that Newton's method yielded significantly faster
convergence rates of \(\cO(\log \log \frac{1}{\epsilon})\). Can we
obtain similar performance for \(\ell_1\)-regularized GLMs?

To obtain a \textbf{proximal Newton method}, we proceed in the same
fashion as above, but rather than approximating the second order term
with \(\alpha_t^{-1} \mbI\), we will use the Hessian of
\(\cL_{\mathsf{d}}\). That leads to a proximal mapping of the form,
\begin{align*}
\mathrm{prox}(\mbu; \mbH_t) 
&= \arg \min_{\mbz} \frac{1}{2} \| \mbz - \mbu \|_{\mbH_t}^2 + \cL_{\mathsf{nd}}(\mbz)
\end{align*} where \(\|\mbx \|_{\mbH_t}^2 = \mbx^\top \mbH_t \mbx\) is a
squared norm induced by the positive definite matrix \(\mbH_t\).

\textit{Note: the proximal mapping for
proximal gradient descent corresponds to the special case in which
\(\mbH_t = \frac{1}{\alpha_t} \mbI\).}

Let \(\mbg_t = \nabla \cL_{\mathsf{d}}(\mbbeta^{(t)})\) and
\(\mbH_t = \nabla^2 \cL_{\mathsf{d}}(\mbbeta^{(t)})\) denote the
gradient and Hessian, respectively. The \emph{undamped} proximal Newton
update is, \begin{align*}
\hat{\mbbeta}^{(t+1)}
&\leftarrow \arg \min_{\mbz} \cL_{\mathsf{d}}(\mbbeta^{(t)}) + (\mbz - \mbbeta^{(t)})^\top \mbg_t + \frac{1}{2} (\mbz - \mbbeta^{(t)})^\top \mbH_t (\mbz - \mbbeta^{(t)}) + \cL_{\mathsf{nd}}(z) \\
&= \arg \min_{\mbz} \frac{1}{2} \| \mbz - (\mbbeta^{(t)} -  \mbH_t^{-1} \mbg_t)\|_{\mbH_t}^2 + \cL_{\mathsf{nd}}(\mbz)  \\
&= \mathrm{prox}(\mbbeta^{(t)} -  \mbH_t^{-1} \mbg_t; \mbH_t)
\end{align*}

As with Newton's method, however, we often need to use damped updates,
\begin{align*}
\mbbeta^{(t+1)} &= \mbbeta^{(t)} + \alpha_t (\hat{\mbbeta}^{(t+1)} - \mbbeta^{(t)}),
\end{align*} 
However, solving the proximal Newton
mapping can be more challenging.

\end{frame}

\begin{frame}[t,allowframebreaks]{Proximal Newton for Sparse
GLMs}\label{proximal-newton-for-sparse-glms}

Let's consider the proximal Newton mapping for \(\ell_1\)-regularized
GLMs, like logistic regression. Here, the non-differentiable part of the
objective is \(\cL_{\mathsf{nd}}(\mbbeta) = \lambda \|\mbbeta\|_1\).
Unfortunately, the proximal Newton update no longer has a closed form
solution because when we introduce the Hessian, the problem no longer
separates across coordinates since the Hessian is generally not
diagonal.

However, note that the proximal Newton step minimizes a second-order
Taylor approximation of the log likelihood plus an
\(\ell_1\)-regularization penalty, \begin{align*}
\mathrm{prox}(\mbu; \mbH_t) 
&= \arg \min_{\mbz} \cL_{\mathsf{d}}(\mbbeta^{(t)}) + (\mbz - \mbbeta^{(t)})^\top \mbg_t + \frac{1}{2} (\mbz - \mbbeta^{(t)})^\top \mbH_t (\mbz - \mbbeta^{(t)}) + \cL_{\mathsf{nd}}(z)
\end{align*}

This is exactly the same problem in the inner loop of the
\texttt{glmnet} algorithm! In particular, we can view the second-order
Taylor approximation of the log likelihood as a weighted least squares
objective with working responses and weights. We can solve that inner
problem with coordinate ascent, just like above.

\textbf{TL;DR: the intuitive algorithm we derived above is really a
proximal Newton algorithm.}
\end{frame}

\begin{frame}[t,allowframebreaks]{Caveats}\label{caveats}

As with regular Newton's method, proximal Newton exhibits local
quadratic convergence to obtain error \(\epsilon\) in
\(\cO(\log \log 1/\epsilon)\) iterations. Though here, each iteration
requires an inner coordinate descent loop to solve the proximal mapping.

\textbf{Warning:} In practice, you may need to also
implement a backtracking line search to choose the step size
\(\alpha_t\), since you may not start in the local quadratic regime.
Logistic regression with decent initialization is reasonably well
behaved, but Poisson regression with log link functions can be
sensitive. 
\end{frame}

\begin{frame}[t,allowframebreaks]{Conclusion}\label{conclusion}

The proximal methods disussed today are what run behind the scenes of
modern packages for sparse linear and logistic regression. 

In
particular,
\href{https://scikit-learn.org/stable/modules/generated/sklearn.linear_model.Lasso.html}{\texttt{sklearn.linear\_model.Lasso}}
uses a fast coordinate descent algorithm like discussed above, and
\href{https://glmnet.stanford.edu/articles/glmnet.html}{GLMNet}
(Friedmann et al., 2010) uses a proximal Newton
algorithm with coordinate descent for the proximal step.

\end{frame}

    
    
\end{document}
